% =========================================================================
% CAMERA-READY REPLACEMENT SNIPPETS  --  drafted from the logs, no re-runs
% =========================================================================
%
% Every number below was recomputed from llm_task_logs/*.json,
% baseline_logs/baseline_results.csv and experiment_logs/experiment_results.csv
% over the frozen window (runs on/after 2026-06-25). Wilson 95% intervals
% recomputed to match.
%
% Replaces: Table V (tab:results), the S1/S2/S3 result subsections, and adds
% the S4 subsection. Also adds a data-provenance paragraph to Sec. V.
%
% CHANGES OF SUBSTANCE vs. the accepted version:
%   1. The single "S1, 16 trials, 81%" row is split into the three scenarios
%      it actually pools (S1 nominal, S2 grounding, S4 disturbance). The
%      pooled figure is kept as a summary line so the accepted number is
%      still visible and reconcilable, not silently withdrawn.
%   2. S4 is reported. It was run; the accepted text says it was not.
%   3. Two-cube ordered delivery is 5/7, not 2/4 -- three further two-cube
%      runs on 2026-07-01 (including one with the order reversed) were
%      omitted from the accepted version.
%   4. The three-cube command ("move all cubes...", 0/9) is reported rather
%      than dropped. It is the evidence for the scalability limitation.
%   5. The baseline row states which implementation produced it.
%   6. S4 is reported descriptively, NOT as an adaptability result: the pick
%      tool resolves coordinates from live perception, so the displacement is
%      absorbed by the architecture, not by the planner. No S4 run called
%      detect_objects, and 2 of the 4 end in place/nav timeouts.
%   7. New subsection sec:verif accounts for self-reported vs verified
%      success and bounds the single-object rate at 9/16--13/16.
%
% =========================================================================


% ---------------------------------------------------------------- Sec. V
% Data provenance -- insert at the end of Sec. V (Experimental Setup),
% before the Results section.

\subsection{Data Provenance}
\label{sec:provenance}

Two logging paths exist in the release and they are not interchangeable, so
we state which produced each reported number. Per-task LLM logs
(\texttt{llm\_task\_logs/*.json}) are written by the planner itself and
record the full conversation, tool calls, timings and token counts; the
experiment runner additionally appends one summary row per trial to
\texttt{experiment\_logs/experiment\_results.csv}. The LLM results below are
computed from the per-task logs, and the scenario each run belongs to is
recovered from the runner's row for the same run. Ten single-object and
two-cube runs on 2026-07-01/02 were launched directly against the planner's
action interface rather than through the runner; these have per-task logs but
no runner row, and since the mid-task disturbance is applied only by the
runner, they are nominal-condition runs by construction. The baseline results
are from \texttt{baseline\_logs/baseline\_results.csv}, written by
\texttt{baseline\_controller}.

The runner also contains a second, inlined copy of the baseline sequence,
used when it is invoked with \texttt{--system baseline}. That copy applies a
0.15\,m approach nudge, whereas \texttt{baseline\_controller} applies the
calibrated 0.30\,m nudge (Sec.~\ref{sec:lessons}); the two are therefore
different controllers and we do not pool their results. Runs produced by the
inlined copy are reported separately in Table~\ref{tab:runner-baseline} and
are not used for any comparison against the LLM planner.


% ---------------------------------------------------------- Table V (new)
% Replaces tab:results.

\begin{table*}[t]
\caption{Frozen-System Results (runs on/after 2026-06-25). Success is as
reported by the system under test and is \emph{not} verified against the
simulator; four of the LLM successes end with a failed or timed-out
\texttt{place} call (Sec.~\ref{sec:verif}). LLM successes typically include
$\geq 1$ autonomous replan.}
\label{tab:results}
\centering
\begin{tabular}{@{}llccccc@{}}
\toprule
\textbf{Scenario} & \textbf{System} & \textbf{N} & \textbf{Success} &
\textbf{95\% CI$^\ddagger$} & \textbf{Mean time (success)} & \textbf{Mean calls} \\ \midrule
S1 single, nominal        & LLM      & 5  & 80\% (4/5)    & [38, 96]\%  & 430\,s (243--714) & 10.2 \\
S2 single, colour grounding & LLM    & 7  & 71\% (5/7)    & [36, 92]\%  & 471\,s (232--714) & 9.0 \\
S4 single, mid-task disturbance & LLM & 4 & 100\% (4/4)   & [51, 100]\% & 457\,s (299--546) & 6.5 \\
\quad\emph{all single-object pooled} & LLM & 16 & \emph{81\% (13/16)} & \emph{[57, 93]\%} & \emph{454\,s (232--714)} & \emph{8.8} \\ \midrule
S1/S2 single              & Baseline$^{\dagger}$ & 6 & 83\% (5/6) & [44, 97]\% & 194\,s (143--249) & --- \\ \midrule
S3 two-cube, ordered      & LLM      & 7  & 71\% (5/7)    & [36, 92]\%  & 551\,s (477--648) & 19.7 \\
S3$'$ three-cube, ordered & LLM      & 9  & 0\% (0/9)     & [0, 30]\%   & ---               & 24.4 \\
S3 two-cube, ordered      & Baseline$^{\dagger}$ & 1 & 0\% (0/1) & [0, 79]\% & --- & --- \\ \midrule
S5 ``tidy up''            & LLM      & 1  & qualitative   & ---         & ---               & 28.0 \\ \bottomrule
\end{tabular}
\vspace{2pt}
\begin{flushleft}
\footnotesize
$^{\dagger}$\texttt{baseline\_controller} on its calibrated 0.30\,m approach
nudge. A seventh single-object run that same session deliberately re-tested
the superseded 0.15\,m nudge and failed at the grasp; it is excluded as a
configuration probe rather than a trial. See
Table~\ref{tab:runner-baseline} for the runner's separate inlined baseline.
$^{\ddagger}$Wilson score interval. Sample sizes are small and all rates are
\emph{preliminary}; the intervals overlap almost everywhere, so we draw
qualitative rather than quantitative conclusions from the comparison.
\end{flushleft}
\end{table*}


% ------------------------------------------------- Table (new, secondary)

\begin{table}[t]
\caption{Trials from the Experiment Runner's Inlined Baseline Sequence
(0.15\,m approach nudge). Reported for completeness; not comparable to
Table~\ref{tab:results}.}
\label{tab:runner-baseline}
\centering
\begin{tabular}{@{}lccl@{}}
\toprule
\textbf{Scenario} & \textbf{N} & \textbf{Success} & \textbf{Dominant failure} \\ \midrule
S1 & 13 & 31\% (4/13) & Navigation (7 of 9)$^{\S}$; grasp reach (2 of 9) \\
S2 & 6  & 17\% (1/6)  & Navigation (5 of 5) \\
S4 & 5  & 20\% (1/5)  & Navigation (3 of 4); grasp reach (1 of 4) \\ \bottomrule
\end{tabular}
\vspace{2pt}
\begin{flushleft}\footnotesize
$^{\S}$Four Nav2 timeouts or failures and three aborts in which the
\texttt{navigate\_to\_pose} action server was not available at trial start.
\end{flushleft}
\end{table}


% ------------------------------------------------------- S4 subsection

\subsection{S4: Behaviour Under a Mid-Task Disturbance}
\label{sec:s4}

S4 repeats the S1 delivery, but the runner teleports the red cube to a new
tabletop position 5\,s after \texttt{navigate\_to(table)} is issued --- while
the robot is still driving. The displacement is $\approx$18\,cm in $y$,
chosen to stay inside the tabletop's collision bounds (a teleport sets only
the instantaneous pose; physics resumes immediately, so a cube placed past
the edge falls) and away from the far edge, where reachability rather than
replanning would dominate the outcome.

All four frozen-system trials were reported complete, but the mechanism is
not the one the disturbance was designed to probe, and we state that
plainly. The \texttt{pick} tool takes an object label, not coordinates: the
tool server resolves the target from live perception at call time. The
planner therefore never holds a stale position to be invalidated. In two of
the four runs the post-disturbance \texttt{pick} simply succeeded at the new
location with no corrective action by the planner at all; in the other two
the planner executed the same forward-nudge recovery it uses in S1, in
response to an out-of-reach error rather than to the displacement.
\emph{No S4 run issued \texttt{detect\_objects}.} What S4 demonstrates is
that a perception-resolved action interface absorbs object displacement ---
an architectural property, available to any controller built on the same
tools --- and not that the language model replans around it.

Two of the four runs are additionally weak as evidence of delivery: the
final \texttt{navigate\_to} and \texttt{place} calls both returned the
planner's 200\,s client-side timeout, after which the model called
\texttt{task\_complete(success=true)} regardless (Sec.~\ref{sec:verif}).
We therefore report S4 as a descriptive result and draw no reliability or
adaptability claim from it. A disturbance experiment that would support such
a claim needs a target the planner must re-localize itself, a
configuration-matched baseline, and the ground-truth delivery check of
Sec.~\ref{sec:verif}; we identify it as the immediate next experiment rather
than claiming it here.


% ------------------------------------------------ S3 subsection (revised)

\subsection{S3: Multi-Object Sequential Delivery}
\label{sec:s3}

Across seven two-cube ordered deliveries the planner succeeded in five
(71\%, [36, 92]\%), averaging 19.7 model calls. Four used the phrasing
``move the red cube to dispatch, then the green cube'' and three a more
explicit variant; one of the latter reversed the order (green first, then
red) and was executed in the stated order, which is evidence that the
sequence is read from the command rather than fixed by the tool set. The two
failures both terminated at the turn limit after repeated replans.

Reliability does not survive a third object. On nine trials of ``move all
cubes to the dispatch zone, red one first,'' the planner never completed the
task (0/9, [0, 30]\%), averaging 24.4 model calls; a further five runner
trials of the same command aborted before the planner produced any call and
are excluded. The characteristic trajectory is that the first cube is
delivered, after which repeated navigation and pick failures accumulate and
the planner exhausts its turn budget while trying to recover. We report this
as the paper's clearest scalability limitation: the same feedback loop that
supplies recovery on a single object supplies, on a longer horizon, an
unbounded supply of things to recover from.


% ------------------------------------------- Success accounting (new)

\subsection{How Success Was Determined, and Its Limits}
\label{sec:verif}

Success in Table~\ref{tab:results} is as reported by the system under test:
for the LLM, its own \texttt{task\_complete(success)} call; for the
baseline, the result of its final step. Neither is checked against the
simulator. Two properties of the stack make this a real limitation rather
than a formality. First, \texttt{place} returns success unconditionally,
including on the fallback path where pre-place IK fails and the object is
released from wherever the arm happens to be. Second, the planner's per-tool
timeout is client-side: when it expires the planner records a failure while
the underlying service call may still be running to completion, so a
timed-out step is genuinely ambiguous rather than definitely failed.

Auditing the tool traces of the 13 self-reported single-object successes
against their final \texttt{place} call: nine end with \texttt{place}
reporting success; three end with a 200\,s \texttt{place} timeout whose true
outcome is unknown; and one reports \texttt{``No object currently held''},
which cannot correspond to a delivery. Read strictly, the single-object rate
is bounded below by 9/16 (56\%) and above by the 13/16 (81\%) we report;
the S4 row is the most affected, with two of its four runs in the
ambiguous class.

We have since added a simulator-side check that reads the cube's true pose
and tests it against the dispatch-zone geometry, recorded as a separate
column alongside the self-reported outcome so the two can disagree visibly.
It cannot be applied retroactively to the runs reported here, and all
future trials in this system are verified by it.
